\documentclass[a4 paper]{article}
% Set target color model to RGB
\usepackage[inner=2.0cm,outer=2.0cm,top=2.5cm,bottom=2.5cm]{geometry}
\usepackage{setspace}
\usepackage[rgb]{xcolor}
\usepackage{verbatim}
\usepackage{subcaption}
\usepackage{amsgen,amsmath,amstext,amsbsy,amsopn,tikz,amssymb}
\usepackage{fancyhdr}
\usepackage[colorlinks=true, urlcolor=blue,  linkcolor=blue, citecolor=blue]{hyperref}
\usepackage[colorinlistoftodos]{todonotes}
\usepackage{rotating}
\usepackage{enumitem}
%\usetikzlibrary{through,backgrounds}
\hypersetup{%
pdfauthor={Karthik Narasimhan},%
pdftitle={Homework},%
pdfkeywords={Tikz,latex,bootstrap,uncertaintes},%
pdfcreator={PDFLaTeX},%
pdfproducer={PDFLaTeX},%
}
%\usetikzlibrary{shadows}
% \usepackage[francais]{babel}
\usepackage{booktabs}
% \input{macros.tex}

\newcommand{\ra}[1]{\renewcommand{\arraystretch}{#1}}

\newtheorem{thm}{Theorem}[section]
\newtheorem{prop}[thm]{Proposition}
\newtheorem{lem}[thm]{Lemma}
\newtheorem{cor}[thm]{Corollary}
\newtheorem{defn}[thm]{Definition}
\newtheorem{rem}[thm]{Remark}
\numberwithin{equation}{section}

\setlength{\parindent}{0pt}

\newcommand{\homework}[6]{
   \pagestyle{myheadings}
   \thispagestyle{plain}
   \newpage
   \setcounter{page}{1}
   \noindent
   \begin{center}
   \framebox{
      \vbox{\vspace{2mm}
    \hbox to 6.28in { {\bf COS 484:~Natural Language Processing \hfill {\small (#2)}} }
       \vspace{6mm}
       \hbox to 6.28in { {\Large \hfill #1  \hfill} }
       \vspace{6mm}
       \hbox to 6.28in {  {\it Name: {\rm #3},
        NetID: {\rm #4} \hfill Students discussed with: {\rm #5}} }
      \vspace{2mm}}
   }
   \end{center}
   \markboth{#3 -- #1}{#3 -- #1}
   \vspace*{4mm}
}

\newcommand{\problem}[1]{~\\\fbox{\textbf{Problem #1}}}
\newcommand{\subproblem}[1]{~\newline\textbf{#1}\newline}
\newcommand{\D}{\mathcal{D}}
\newcommand{\Hy}{\mathcal{H}}
\newcommand{\VS}{\textrm{VS}}

\newcommand{\bbF}{\mathbb{F}}
\newcommand{\bbX}{\mathbb{X}}
\newcommand{\bI}{\mathbf{I}}
\newcommand{\bX}{\mathbf{X}}
\newcommand{\bY}{\mathbf{Y}}
\newcommand{\bepsilon}{\boldsymbol{\epsilon}}
\newcommand{\balpha}{\boldsymbol{\alpha}}
\newcommand{\bbeta}{\boldsymbol{\beta}}
\newcommand{\0}{\mathbf{0}}
\newcommand{\code}[1]{\texttt{#1}}


\begin{document}
\homework{Assignment \#3}{Due: 03/31/2025}{Athan Zhang}{az7019}{None}
\textbf{Course Policy}: Read all the instructions below carefully before you start working on the assignment, and before you make a submission. The course assignment policy is available at \url{http://nlp.cs.princeton.edu/cos484}.

\newpage
\section*{Problem 1: Language modeling with recurrent neural networks}

\subsection*{Part A}

The parameters of this neural network that can be learned from the training data are:

\begin{itemize}
    \item The input embedding $E$: $|V| \times D$ parameters
    \item The input weight matrix $W$: $D\times H$ parameters
    \item The recurrent weight matrix $U$: $H \times H$ parameters
    \item The output embedding $W_O$: $|V| \times H$ parameters
\end{itemize}

Thus the total number of trainable parameters the RNN has is 

\begin{align*}
    \boxed{|V|D + |V|H + H^2 + DH}
\end{align*}

\subsection*{Part B}

\begin{align*}
    \frac{\partial y_3}{\partial E} &= \frac{\partial y_3}{\partial h_3} \frac{\partial h_3}{\partial E} \\
    &= \frac{\partial y_3}{\partial h_3} \left( \frac{\partial h_3}{\partial e_3} \frac{\partial e_3}{\partial E} + \frac{\partial h_3}{\partial h_2} \frac{\partial h_2}{\partial E} \right) \\
    &= \frac{\partial y_3}{\partial h_3} \left( \frac{\partial h_3}{\partial e_3} \frac{\partial e_3}{\partial E} + \frac{\partial h_3}{\partial h_2} \left( \frac{\partial h_2}{\partial e_2} \frac{\partial e_2}{\partial E} + \frac{\partial h_2}{\partial h_1} \frac{\partial h_1}{\partial E} \right) \right) \\
    &= \boxed{\frac{\partial y_3}{\partial h_3}\frac{\partial h_3}{\partial e_3} \frac{\partial e_3}{\partial E} + \frac{\partial y_3}{\partial h_3}\frac{\partial h_3}{\partial h_2}\frac{\partial h_2}{\partial e_2}\frac{\partial e_2}{\partial E} + \frac{\partial y_3}{\partial h_3}\frac{\partial h_3}{\partial h_2} \frac{\partial h_2}{\partial h_1} \frac{\partial h_1}{\partial e_1} \frac{\partial e_1}{\partial E}}
\end{align*} 

\subsection*{Part C}

One reason why weight tying might be beneficial is that the model learns a consistent embedding dimension for words both when they appear/are parsed for input and when they are predicted for output. This can lead to improved performance and less overfitting.

\subsection*{Part D}

If $U$ is a $2 \times 2$, to have the terms tend to 0, we want the values on the diagonal to be less than 1. An example would be 

\begin{align*}
    U = \boxed{\begin{bmatrix}
        0.5 & 0 \\ 0 & 0.5
    \end{bmatrix}}
\end{align*}

For terms to tend toward infinity, we want the values of the diagonal to be greater than 1. An example would be

\begin{align*}
    U = \boxed{\begin{bmatrix}
        5 & 0 \\ 0 & 5
    \end{bmatrix}}
\end{align*}

\newpage
\section*{Problem 2: IBM Models}

\subsection*{Part A}

With a uniform translation model, we know that all translations occur with the exact same probability. 

\begin{align*}
    P(w^{(s)} | w^{(t)}) = \frac{1}{|V^t|} = \boxed{\frac{1}{4}}
\end{align*}

\subsection*{Part B}

Each target-sentene word may align to any of the source-sentence words. Thus, if the source has $m$ words and the target has $n$ words, then there are $m^n$ ways to choose one-to-one alignments.

For the first sentence, we have $3^3 = 27$ possible alignments and for the second sentence we have $2^2 = 4$ possible alignments.

\subsection*{Part C}

The only word pairs that will have nonzero probabilities will be words that are represented somewhere from the given alignment. Felines is aligned with cats and hats once each. Thus, felines has two total alignments, and each case occurred once. Thus, we have the following probabilities:

\begin{enumerate}
    \item $P(cats | felines) = \frac{1}{2}$
    \item $P(hats | felines) = \frac{1}{2}$
    \item $P(kids | children) = 1$
    \item $P(like | adore) = 1$
    \item $P(cats | fedora) = 1$
\end{enumerate}

\subsection*{Part D}

The most likely translation of the new sentence "cats like hats" is identified by finding the value of $w^{(t)}$ that maximizes $P(w^{(t)} | w^{(s)})$. 

\begin{align*}
    P(w^{(t)} | w^{(s)}) \propto  P(w^{(t)})\prod_{j}^{M^{(s)}}\sum_{i}^{M^{(t)}} P(w_{j}^{(s)} | w_{i}^{(t)})
\end{align*}

Here are the following evaluations for all translations:
\newline 

"fedoras adore fedoras":
\begin{align*}
    &P(\text{fedoras adore fedoras}) \times \\
    &\quad \Big[ 
        \left( P(\text{cats} \mid \text{fedoras}) + P(\text{cats} \mid \text{adore}) + P(\text{cats} \mid \text{fedoras}) \right) \times \\
    &\quad \left( P(\text{like} \mid \text{fedoras}) + P(\text{like} \mid \text{adore}) + P(\text{like} \mid \text{fedoras}) \right) \times \\
    &\quad \left( P(\text{hats} \mid \text{fedoras}) + P(\text{hats} \mid \text{adore}) + P(\text{hats} \mid \text{fedoras}) \right) 
    \Big] \\
    &= 0.2 \cdot \left[ (1 + 0 + 1)(0 + 1 + 0)(0 + 0 + 0) \right] \\
    &= (0.2)(2)(1)(0) \\
    &= 0
\end{align*}
\newline

"felines adore fedoras":
\begin{align*}
    &P(\text{felines adore fedoras}) \times \\
    &\quad \Big[ 
        \left( P(\text{cats} \mid \text{felines}) + P(\text{cats} \mid \text{adore}) + P(\text{cats} \mid \text{fedoras}) \right) \times \\
    &\quad \left( P(\text{like} \mid \text{felines}) + P(\text{like} \mid \text{adore}) + P(\text{like} \mid \text{fedoras}) \right) \times \\
    &\quad \left( P(\text{hats} \mid \text{felines}) + P(\text{hats} \mid \text{adore}) + P(\text{hats} \mid \text{fedoras}) \right) 
    \Big] \\
    &= 0.25 \cdot \left[ (0.5 + 0 + 1)(0 + 1 + 0)(0.5 + 0 + 0) \right] \\
    &= (0.25)(1.5)(1)(0.5) \\
    &= 0.1875
\end{align*}
\newline

"children adore felines":
\begin{align*}
    &P(\text{children adore felines}) \times \\
    &\quad \Big[ 
        \left( P(\text{cats} \mid \text{children}) + P(\text{cats} \mid \text{adore}) + P(\text{cats} \mid \text{felines}) \right) \times \\
    &\quad \left( P(\text{like} \mid \text{children}) + P(\text{like} \mid \text{adore}) + P(\text{like} \mid \text{felines}) \right) \times \\
    &\quad \left( P(\text{hats} \mid \text{children}) + P(\text{hats} \mid \text{adore}) + P(\text{hats} \mid \text{felines}) \right) 
    \Big] \\
    &= 0.4 \cdot \left[ (0 + 0 + 0.5)(0 + 1 + 0)(0 + 0 + 0.5) \right] \\
    &= (0.4)(0.5)(1)(0.5) \\
    &= 0.1
\end{align*}
\newline

"adore fedoras felines":
\begin{align*}
    &P(\text{adore fedoras felines}) \times \\
    &\quad \Big[ 
        \left( P(\text{cats} \mid \text{adore}) + P(\text{cats} \mid \text{fedoras}) + P(\text{cats} \mid \text{felines}) \right) \times \\
    &\quad \left( P(\text{like} \mid \text{adore}) + P(\text{like} \mid \text{fedoras}) + P(\text{like} \mid \text{felines}) \right) \times \\
    &\quad \left( P(\text{hats} \mid \text{adore}) + P(\text{hats} \mid \text{fedoras}) + P(\text{hats} \mid \text{felines}) \right) 
    \Big] \\
    &= 0.1 \cdot \left[ (0 + 1 + 0.5)(1 + 0 + 0)(0 + 0 + 0.5) \right] \\
    &= (0.1)(1.5)(1)(0.5) \\
    &= 0.075
\end{align*}
\newline

"adore fedoras fedoras":
\begin{align*}
    &P(\text{adore fedoras fedoras}) \times \\
    &\quad \Big[ 
        \left( P(\text{cats} \mid \text{adore}) + P(\text{cats} \mid \text{fedoras}) + P(\text{cats} \mid \text{fedoras}) \right) \times \\
    &\quad \left( P(\text{like} \mid \text{adore}) + P(\text{like} \mid \text{fedoras}) + P(\text{like} \mid \text{fedoras}) \right) \times \\
    &\quad \left( P(\text{hats} \mid \text{adore}) + P(\text{hats} \mid \text{fedoras}) + P(\text{hats} \mid \text{fedoras}) \right) 
    \Big] \\
    &= 0.05 \cdot \left[ (0 + 1 + 1)(1 + 0 + 0)(0 + 0 + 0) \right] \\
    &= (0.05)(2)(1)(0) \\
    &= 0
\end{align*}

We can see that "felines adore fedoras" is the most likely translation.

\subsection*{Part E}

Let's first calculate the unigram precision. The candidate unigrams are $\{ \text{felines, adore, fedoras} \}$. The reference unigrams are $\{ \text{felines, adore, chapeaus} \}$. We can see there are 2 out of 3 matches.

Next let's calculate the bigram precision. The candidate bigrams are $\{ \text{felines-adore, adore-fedoras} \}$. The reference bigrams are $\{ \text{felines-adore, adore-chapeaus} \}$. We can see there are 1 out of 2 matches.

To calculate our BLEU-2 score, we can use the following formula

\begin{align*}
    BLEU_2 &= \sqrt{prec_1 \times prec_2} \\
    &= \sqrt{\frac{2}{3} \cdot\frac{1}{2}} \\
    &= \boxed{\sqrt{\frac{1}{3}}}
\end{align*}

\end{document}
